\section{Projection-Relative Saturation}

\subsection{Why saturation must be projection-relative}

The finite mechanics contains several distinct forgetful maps.
Consequently, the statement that a system has ``lost information''
is incomplete unless the relevant map is named.

Let
\[
q:X\longrightarrow Y
\]
be a declared projection.

A distinction between \(x,x'\in X\) survives the projection when
\[
q(x)\neq q(x').
\]
It is hidden by the projection when
\[
x\neq x',
\qquad
q(x)=q(x').
\]

The second relation does not imply that the complete states have
merged. It says only that the declared projection no longer resolves
their difference.

\begin{definition}[Projection-relative degeneracy]
Two distinct complete states \(x,x'\in X\) are projection-degenerate
relative to \(q\) when
\[
q(x)=q(x').
\]
\end{definition}

This is an exact statement about a map.

It is not an ontological identification of \(x\) and \(x'\).

\subsection{Independent registered distinction}

Let
\[
\mathcal X\subseteq X
\]
be a declared family of complete states.

The projected state set is
\[
q(\mathcal X)\subseteq Y.
\]

When the relevant notion of independence is cardinal, the number of
projected distinction classes is simply
\[
K_q(\mathcal X)
=
|q(\mathcal X)|.
\]

More generally, when a declared linear or combinatorial rank function
is available, write
\[
r_q(\mathcal X)
\]
for the independent distinction rank that survives the projection.

No universal rank function is assumed in this paper.

The choice must be stated for the particular register being studied.

\subsection{Projection-relative saturation}

\begin{definition}[Projection-relative saturation]
Let
\[
\{x_\lambda\}_{\lambda\in\Lambda}
\]
be a declared loading family of complete states or admitted histories,
and let
\[
q:X\longrightarrow Y
\]
be a declared projection or registration.

The family enters a projection-relative saturation regime when:

\begin{enumerate}[label=(\roman*)]
\item complete distinction remains nontrivial;

\item the independent distinction visible through \(q\) is bounded;

\item increasing loading ceases to produce corresponding growth of
independently resolved projected distinction;

\item additional complete variation is therefore increasingly carried
inside fibres of \(q\), retained receipts, lifted structure, or other
higher-resolution data.
\end{enumerate}
\end{definition}

Saturation is therefore not defined as disappearance of the complete
state.

It is a statement about the relationship between complete variation
and the resolving capacity of a declared registration.

\subsection{Contrast saturation on the sector register}

For the fifteen-sector real register, Paper 14 defines
\[
\zetaRel(x)
=
\frac{\|\Pperp x\|_2^2}{\|x\|_2^2},
\qquad
x\neq0.
\]

This dimensionless quantity measures the fraction of register
magnitude carried by centered relational contrast.

\begin{definition}[Contrast-saturating sequence]
A sequence of nonzero registered states
\[
x_0,x_1,x_2,\ldots
\]
is contrast-saturating when
\[
\zetaRel(x_k)
\]
approaches a stable minimum while the corresponding complete state
sequence remains nontrivial.
\end{definition}

For normalized quadratic aggregation, if
\[
\Pzero x\neq0,
\]
the contrast-concentration theorem gives
\[
\zetaRel(\widehat Q^{\,k}x)
\longrightarrow0.
\]

Thus \(\widehat Q\)-iteration supplies an exact algebraic example of
complete common-mode concentration.

\begin{corollary}[Asymptotic contrast saturation]
If
\[
\Pzero x\neq0,
\]
then
\[
\lim_{k\to\infty}
\zetaRel(\widehat Q^{\,k}x)
=
0.
\]
For every finite \(k\), however,
\[
\widehat Q^{\,k}
\]
remains invertible.
\end{corollary}

Hence the same finite sequence exhibits
\[
\text{relative registered contrast}\longrightarrow0
\]
while
\[
\text{exact finite-stage distinction remains recoverable}.
\]

This is the exact algebraic core motivating the saturation language.

\subsection{Finite-resolution saturation}

Exact arithmetic is not the only possible declared registration.

Let
\[
\varepsilon>0
\]
be a fixed resolution threshold, and let
\[
d=x-y
\]
be a nonzero centered distinction:
\[
d\in\one_{15}^{\perp}.
\]

The contrast-concentration theorem gives
\[
\|\widehat Q^{\,k}d\|_2
\le
\left(\frac{9}{49}\right)^k
\|d\|_2.
\]

Therefore, for every fixed \(\varepsilon>0\), there exists a finite
\(k_\varepsilon\) such that
\[
0<
\|\widehat Q^{\,k}d\|_2
<
\varepsilon
\]
for all
\[
k\ge k_\varepsilon.
\]

The distinction remains mathematically nonzero, but a registration
with resolution \(\varepsilon\) no longer distinguishes it.

We call this
\emph{effective registration degeneracy}.

It must remain distinct from exact equality.

\subsection{Saturation is not global incompatibility}

The conjecture presupposes an admitted globally compatible transport
structure.

Failure of locally admissible choices to assemble into a globally
closed transport relation is a different phenomenon.

We therefore distinguish:

\begin{description}

\item[Local inadmissibility.]
A proposed local transition is not contained in the admitted local
transport vocabulary.

\item[Global incompatibility.]
Locally admissible transitions cannot be assembled into a globally
closed transport assignment satisfying the declared compatibility
conditions.

\item[Relational saturation.]
A globally admitted construction retains complete distinctions while
a declared registration ceases gaining independent resolving power.

\end{description}

The third must never be inferred from the second.

\subsection{Boundary}

The exact result of this section is a definition plus a proved
algebraic example.

Normalized \(Q\)-iteration is contrast-saturating relative to the
centered sector readout.

What remains unproved is that a native admitted Thalean loading law
realizes this sequence or an equivalent contrast-concentrating family.

Accordingly,
\[
\text{contrast saturation under }\widehat Q
\]
is theorem-side, while
\[
\text{native relational saturation}
\]
remains conjectural.
